\documentclass{beamer}
% \usepackage{graphicx}
% \usepackage[all,cmtip]{xy}
\usepackage{tikz}
\usepackage{tikz-cd}
\usepackage{qtree}
% \usepackage{diagxy}
\usepackage{amsfonts}
\usepackage{amssymb}
\usepackage{amsmath}
\usepackage{amscd}
\usepackage{amsthm}
\usepackage{mathpartir}
\usepackage{bbm}
%\usepackage{pictexwd,dcpic}
\usepackage[mathcal]{euscript}
\usepackage[utf8]{inputenc}
%\usepackage[T1]{fontenc}
\usepackage{lmodern}
%\usepackage{setspace}
%\usepackage{mathpartir}
% \usepackage{proof}
%\setstretch{1}
\title{The relativity of wrong, redux}
\author{Hans Halvorson}
\date{May 9, 2013}
\usetheme{Copenhagen}
\begin{document}
\def\2{\mathcal}
\def\7{\mathbb}
\begin{frame}

\maketitle

\end{frame}

\begin{frame}



{\LARGE Is it irrational to believe that God exists?} 

\vspace{2.5em}
{\LARGE Is it irrational to hope for eternal life?}

\vspace{2.5em}
{\LARGE Has science made religion obsolete?} 

\end{frame}

\begin{frame}{The positivist story (Auguste Comte)}

\begin{tabular}{cll}
  Science & &  Newton, Darwin, Hawking  \\ $\Uparrow$ \\
  Philosophy & & Aristotle, Descartes, Kant \\ $\Uparrow$ \\ 
  Religion & & Judaism, Christianity, Islam \\ $\Uparrow$ \\ 
  Myth & & Zeus, Loki  \\  \\
\end{tabular}
\end{frame}

\begin{frame}{Science's victorious march}

\begin{tabular}{lllll}
  & & Religion & Science \\ \hline \hline \\
  2013 & & Soul &  Neuroscience & \\ \\ 
  1860 & & Creation  & Evolution & Darwin \\ \\
  1700 & & Miracles  & Natural Laws & Newton \\ \\
  1600 & & Geocentric & Heliocentric & Galileo \\ \\ 
  1200 & & Creation & Eternal Universe & Aristotle\\ \\ 
  \vdots & & \vdots &  \vdots \end{tabular}
\end{frame}


\begin{frame}{What's wrong with these two stories about science}

\begin{description}
\item[POINT:] Most past scientific theories have been wrong
\item[COUNTERPOINT:] Some theories are wronger than others; science
  develops through ``successive refinements'' (Isaac Asimov: ``The
  relativity of wrong'')
\end{description}

\end{frame}

\begin{frame}{What's wrong with these two stories about science}

\begin{description} 
\item[POINT:] Science is not obviously progressive, not obviously
  converging to the truth \end{description}

\bigskip
\begin{itemize}
\item Wine is good for you (2000).  No, wine is bad for you (2008).
  No, wine is in fact good for you (2013).
\item Personality is due primarily to genetics.  No, it's due
  primarily to environmental conditions.  No, genetics is the most
  important factor.
\end{itemize}



\end{frame}

\begin{frame}{What's wrong with these two stories about science}

\begin{itemize}
\item Science surrenders the territory it apparently gained

\bigskip \begin{tabular}{lll}
eternal universe & $\longrightarrow$ & big bang \\ \\ 
heliocentric & $\longrightarrow$ & no preferred reference frame \\ \\
natural laws & $\longrightarrow$ & quantum mechanics \end{tabular}
\end{itemize}



\end{frame}


\begin{frame}{Case Study: The Eternity of the World}

\begin{itemize}
\item Aristotle: Second great philosopher; first great scientist
\item Rejects Plato's forms.  The world consists of en-formed matter.
\item Rejects Plato's claim that there can be no knowledge of the
  physical world $\longrightarrow$ made physics possible
\end{itemize}
\end{frame}


\begin{frame}{Conflict areas between Aristotle and Religion}

\begin{enumerate}
\item Materialism and the possibility of life after death.
\item Ethics
\item The creation of the universe
\end{enumerate}

\end{frame}

\begin{frame}{Aristotle's Physics}

  \begin{itemize}
  \item In every change, something stays the same (underlying matter).
    What changes is the \emph{form}.
\item Every change is a sort of movement.
\item Nature abhors a vacuum.
\end{itemize}
\end{frame}

\begin{frame}{The Universe is Eternal}

\begin{enumerate} \item Suppose for \emph{reductio ad absurdum} that
  the universe came into existence at some finite time.  Then matter
  existed after that time, but didn't exist before that time.  But the
  change (from non-existence to existence) requires an underlying
  substratum, and so matter existed before the universe existed.  This
  is a contradiction.  
\item Suppose for \emph{reductio ad absurdum} that there was a first
  motion.  Then either the universe existed eternally in a state of
  rest before the first motion, or the universe came into existence
  and then something moved.  The second option is a contradiction,
  because the universe coming into existence would itself be an object
  moving.  The first option is also wrong.
\end{enumerate}
\end{frame}


\begin{frame}{St.\ Bonaventure}

\begin{itemize}
\item Reason was flawed by the fall, and cannot find truth without the
  help of faith.  (Distrust of Aristotle)
\item Reason, when used properly, shows that the universe is finite.
\item Contemporary parallels:
\begin{itemize}
\item Creation Science
\item Intelligent Design
\item Alvin Plantinga: Augustinian Science
\end{itemize}
\end{itemize}

\end{frame}

\begin{frame}{Thomas Aquinas}

\begin{itemize}
\item Criticism of Bonaventure: bad science reflects poorly on
  Christianity (compare Francis Collins on intelligent design)
\item Science cannot demonstrate that the world is eternal, nor that
  it was created.  
\item ``Was the world created?'' Can only be answered by faith
\item Reason cannot contradict faith, but Reason has fixed
  limits. % cf Maimonides
\item Contemporary parallels
\begin{itemize}
\item Still holds sway in many Catholic circles
\item Non-overlapping magesteria 
\end{itemize}
\end{itemize}


\end{frame}
\begin{frame}{Siger of Brabant and Boetius of Dacia}

\begin{itemize}
\item Accused of holding the doctrine of \textbf{double truth}: two
  contradictory things can both be true
  \begin{quote} \emph{As a scientist, I believe that the world is
      eternal.  \\ As a Christian, I believe that the world is not
      eternal.} \end{quote}
\item A scientist can only admit those things that follow the
  principles of his science
\item Science cannot admit creation events
\item As far as science is concerned, every natural event has a
  natural cause
\item Fideism: some things should be believed even though not
  demonstrable (???)
\item Contemporary parallels
\begin{itemize}
\item Methodological naturalism
\end{itemize}
\end{itemize}

\end{frame}


\begin{frame}{If those guys were alive today \dots }

\begin{itemize}
\item Bonaventure: reject mainstream science, develop a Theistic
  alternative
\item Aquinas: show where science leaves off and faith begins
\item Siger and Boetius: science is a serious game
\end{itemize}

\end{frame}

\begin{frame}{The missing point of view}

  None of our medieval thinkers entertains the idea that science could
  answer \emph{all} questions (contrast with Dawkins, Dennett, Harris,
  etc.)

  \bigskip For Aquinas, no strict separation between Reason (Science)
  and Faith.  Aquinas calls Theology a Science.  But distinguish:
\begin{itemize}
\item Things that Reason can prove (e.g.\ God exists).
\item Things that are matters of faith (e.g.\ Jesus is the incarnate
  God, the universe was created in time). \end{itemize}

\end{frame}


\begin{frame}{Scientism: Science can answer all questions}

\begin{itemize}
\item What is science?  Is sociology a science?  Is history a science?
  Is philosophy a science?
  \begin{quote} ``For Quine, philosophical theories are to be assessed
    in the ways that scientific hypotheses are, and good philosophical
    method is continuous with good scientific method.''
\end{quote}
\item Does science tell me what I ate for breakfast this morning?
\item Does science tell me that it's wrong to kill?
\item Does science tell me that a Brahms concerto is beautiful?
\end{itemize}


\end{frame}


\begin{frame}{The Galileo Affair}

\begin{itemize}
\item Christian Doctrine before 1600: the earth is the center of the
  universe (geocentrism)
\item Geocentrism fits the data only at the price of horrendous
  complications (epicycles)
\item The Copernican revolution: everything becomes simple if we put
  the sun in the center
\item Galileo condemned for teaching heliocentrism
\end{itemize}

\end{frame}

\begin{frame}{They were all wrong}

\begin{itemize}
\item Does the earth go around the sun, or the sun around the earth?
\item Galilean relativity: No such thing as absolute rest
\item General relativity: No such thing as absolute rotation \\
\smallskip  \emph{There is no objective right answer about whether $A$ is
    rotating around $B$ or vice versa}.
\item Lesson: The debate ($X$ or not-$X$) was based on a false
  presupposition.
\end{itemize}

\end{frame}



\begin{frame}{Prospective views on science and religion}

\begin{itemize}
\item Secular scientism: Science and religion are in competition, and
  science wins
\item Religious scientism: Science will eventually come around to the
  old truths of religion
\item Separatism: Science and religion are non-overlapping magesteria
  (Steven Jay Gould).  Science has to do with facts, religion with
  values.
\item A~science is a ``rule governed activity'' with limited
  objectives and limited possibilities.  \begin{quote} \emph{According
      to the principles of this science\dots } \end{quote}
\end{itemize}


\end{frame}


\end{document}



\begin{frame}

  Question 1: Is it irrational to believe that God exists, and that
  God has revealed Himself to us?

  \bigskip Question 2: If a person believes $X$, and then they
  discover evidence that supports not-$X$, would it be irrational to
  continue believing $X$?

\bigskip In cases of disagreement, some general answers:
\begin{itemize}
\item Split the difference (Elga)
\item Stick to your guns (Kelly)
\end{itemize}

\end{frame}





\begin{frame}

  Specific Focus: (a) $X$ is some sort of religious belief, and (b)
  there is scientific evidence against $X$.

  \bigskip Examples:
\begin{itemize}
\item $X$ = God created the universe in six days; the first man was
  Adam, and his sin brought suffering and death into the world.
\item $X$ = We have the freedom to choose whether to do good or to do
  bad.
\item $X$ = God created the universe.
\end{itemize}

\end{frame}

\begin{frame}{General Strategies for Responding to Contrary Scientific
    Data}

\begin{enumerate}
\item Hold on to the religious belief and reject the new scientific
  data.
\item Accept the new scientific data and reject the religious beliefs.
\item Try to hold on to both religion and science.
\end{enumerate}

\end{frame}

